\documentclass[10pt,twocolumn]{article}

\usepackage[margin=0.72in]{geometry}
\usepackage{times}
\usepackage{url}
\usepackage{graphicx}

\title{Style-Aware Adaptation for Language-Conditioned Humanoid Motion Generation}
\author{Haotian Wu \and Yuhua Luo \and Xuyang Yuan}
\date{}

\begin{document}
\maketitle

\begin{abstract}
This midterm report summarizes progress toward style-aware adaptation of MoConVQ for language-conditioned humanoid motion generation. We reproduced the pretrained MoConVQ inference pipeline, generated ten baseline BVH motions, built split-preserving HumanML3D style-subset tooling, and identified a key practical constraint: the available HumanML3D parquet export contains processed motion features rather than MoConVQ-compatible token targets. The next phase therefore uses parameter-efficient LoRA distillation as a faithful backup to the proposal's frozen-backbone adaptation plan.
\end{abstract}

\section{Problem and Motivation}
The goal is to generate humanoid motion from natural language while preserving both action semantics and requested style. MoConVQ is a strong robotics baseline because it uses discrete motion representations and physics-based tracking control \cite{yao2024moconvq}. However, style-specific prompts such as slow deliberate walking, tai chi-like movement, graceful dance, and heavy tired walking may be underrepresented in general text-to-motion training data.

\section{Related Work}
HumanML3D provides paired motion and text descriptions for text-driven motion generation \cite{guo2022humanml3d}. MDM models raw motion with diffusion \cite{tevet2023mdm}, while T2M-GPT demonstrates the effectiveness of discrete motion tokens for text-to-motion generation \cite{zhang2023t2mgpt}. Our project follows the discrete-token direction but keeps MoConVQ's motion representation and controller frozen.

\section{Technical Approach}
The intended proposal pipeline freezes the VQ representation, codebook, decoder, and controller, while adapting only the language-conditioned code generator.

\begin{figure}[h]
\centering
\fbox{\begin{minipage}{0.94\linewidth}
\centering
Text prompt $\rightarrow$ text encoder $\rightarrow$ language transformer $\rightarrow$ discrete motion codes $\rightarrow$ MoConVQ decoder/controller $\rightarrow$ BVH humanoid motion
\end{minipage}}
\caption{Project pipeline. Adaptation is restricted to the language-conditioned generator; the motion backbone remains frozen.}
\end{figure}

For data curation, we filter style-related captions independently within each HumanML3D split using five keyword groups: controlled motion, graceful smooth motion, expressive dance, energetic dynamics, and heavy/awkward motion. This prevents leakage across train, validation, and test sets.

\section{Midterm Progress}
Baseline reproduction is complete. The official pretrained MoConVQ assets are present locally, tokenization and decoding have been verified, and ten prompts have produced BVH files under \texttt{outputs/baseline/project\_prompts/}.

\begin{table}[h]
\centering
\caption{Implementation progress at midterm.}
\footnotesize
\begin{tabular}{p{0.30\linewidth}p{0.17\linewidth}p{0.31\linewidth}}
\hline
Component & Status & Evidence \\
\hline
Environment & Complete & CUDA PyTorch and local extensions verified \\
Baseline inference & Complete & 10 BVH files generated \\
Style filtering & Complete & Split-preserving parquet filter \\
Small splits & Complete & 200/40/40 style subset files \\
Fine-tuning & Planned & LoRA distillation selected after format diagnosis \\
\hline
\end{tabular}
\end{table}

\begin{table}[h]
\centering
\caption{Style subset statistics from the available HumanML3D parquet export.}
\footnotesize
\begin{tabular}{lrr}
\hline
Split & Matched captions & Matched motion IDs \\
\hline
Train & 3332 & 2682 \\
Validation & 214 & 168 \\
Test & 586 & 486 \\
\hline
\end{tabular}
\end{table}

\section{Risk and Next Step}
The main technical risk is format mismatch. The available HumanML3D parquet export contains processed 263-D features, while MoConVQ supervised fine-tuning requires discrete token sequences in its own representation. A direct conversion is nontrivial and could introduce invalid supervision.

The final-stage plan is therefore to use a parameter-efficient backup: pretrained MoConVQ serves as a teacher to generate pseudo token targets for style captions, and a student trains only LoRA updates in the language-conditioned transformer. We will report only measured loss, generated BVH outputs, and qualitative comparisons unless a compatible FID/R-Precision evaluator is implemented.

\bibliographystyle{plain}
\bibliography{../shared/references}

\end{document}
